\documentclass[11pt,a4paper]{article}
% CM3191 Organic Experiment 2 -- report draft (pre-written). Compile with XeLaTeX.
\newif\ifdraft \drafttrue
% ---------- common preamble : CM3191 organic report drafts (XeLaTeX) ----------
\usepackage[a4paper,margin=22mm]{geometry}
\usepackage{fontspec}
\usepackage{xeCJK}
\setmainfont{TeX Gyre Termes}
\setsansfont{TeX Gyre Heros}
\setCJKmainfont{Noto Serif CJK KR}
\xeCJKsetup{CJKspace=true, CJKecglue={}}
\usepackage{amsmath,amssymb}
\usepackage[version=4]{mhchem}
\usepackage{chemfig}
\usepackage{siunitx}
\sisetup{detect-all, per-mode=symbol}
\usepackage{booktabs,tabularx,array}
\usepackage{enumitem}
\usepackage{xcolor}
\usepackage[super]{cite}
\usepackage{titlesec}
\usepackage[hidelinks]{hyperref}
\usepackage{caption}
\usepackage{float}
\captionsetup{font=small,labelfont=bf}
\titleformat{\section}{\large\bfseries}{\thesection.}{0.5em}{}
\titleformat{\subsection}{\normalsize\bfseries}{\thesubsection}{0.5em}{}
\setlength{\parskip}{3pt}
\setlength{\emergencystretch}{3em}
\sloppy
\linespread{1.1}
\setchemfig{atom sep=1.9em, bond style={line width=0.6pt}}

% ---- placeholders: everything in a grey box must be replaced after the lab
\newcommand{\res}[1]{\fcolorbox{black}{gray!18}{\footnotesize\sffamily #1}}
\newcommand{\obs}[1]{\fcolorbox{black}{gray!18}{\footnotesize\sffamily OBS: #1}}
\newcommand{\resp}[1]{\par\smallskip\noindent\fcolorbox{black}{gray!18}{\parbox{\dimexpr\linewidth-2\fboxsep-2\fboxrule\relax}{\footnotesize\sffamily #1}}\par\smallskip}

% ---- title block: \reporttitle{title}{experiment line}{date}
\newcommand{\reporttitle}[3]{%
\begin{center}
{\LARGE\bfseries #1}\\[4mm]
{\large Kim Donghwan (A0266094X)}\\[1mm]
CM3191 Chemical Synthesis Experiments, #2\\
Thursday laboratory, fume hood B7 \quad\textbullet\quad Date of experiment: #3\\
Instructor: Dr Hoang Truong Giang
\end{center}
\vspace{2mm}
\noindent\rule{\linewidth}{0.6pt}}

% ---- Korean drafting notes page (removed with \draftfalse)
\newcommand{\draftnotes}[2]{%
\ifdraft
\thispagestyle{empty}
\begin{center}{\Large\bfseries 작성 안내 (제출 전 이 페이지 삭제)}\end{center}
\vspace{2mm}
\textbf{이 초안의 구조.} 결과와 무관한 부분(서론·이론, 실험 방법, 고찰의 이론적 틀, 참고문헌)은 완성된 문장으로 써 두었습니다. 실험 후 바뀌는 값과 관찰은 \res{회색 상자}로 남겼습니다. 상자 안의 숫자는 매뉴얼의 계획값이므로 \textbf{반드시 본인 측정값으로 바꿔야 합니다}. 결과에 따라 문장이 달라지는 곳은 .tex 파일에 \texttt{\%} 주석으로 선택지를 달아 두었습니다.

\textbf{실험 후 채울 것}
\begin{enumerate}[leftmargin=6mm,itemsep=0.5mm]
#1
\end{enumerate}

\textbf{제출 전 확인할 것 [CONFIRM]}
\begin{enumerate}[leftmargin=6mm,itemsep=0.5mm]
\item 유기 파트 보고서 형식·분량·마감(평가 방식 공지 확인). proforma 형식이면 문단을 답안으로 옮겨 쓰면 됩니다.
\item 과목의 생성형 AI 사용 규정(사용 사실 표기 필요 여부).
#2
\end{enumerate}
\textbf{컴파일}: XeLaTeX. 파일 첫 줄 근처의 \texttt{\textbackslash drafttrue}를 \texttt{\textbackslash draftfalse}로 바꾸면 이 페이지가 빠집니다.
\newpage
\setcounter{page}{1}
\fi}


\begin{document}

\draftnotes{%
\item 2절과 표 1: 실제 질량, KOH 방울 수, 추가 styrene 여부, 재결정에 더 넣은 에틸아세테이트 양.
\item 3절: TLC 거리와 $R_f$(표 3, datasheet 그림 옮기기), 생성물 질량·수율, 녹는점, IR, \textsuperscript{1}H NMR의 비방향족 H 세 개(δ, $J$), X선 구조로 확정한 생성물 번호.
\item 수율 계산 기준: X선 결과가 \textbf{3A/3B}이면 MW 488.63, \textbf{4}이면 460.62. 표 2와 초록의 해당 줄만 남깁니다.
\item 4.3--4.4절: 주석의 선택지 중 결과에 맞는 문장 선택.}{%
\item NMR 측정 주파수(400 MHz로 가정), IR 기기명, 녹는점 측정기.
\item 강의에서 제공될 X선 구조 자료를 그림으로 넣을지.}

\reporttitle{One-Pot Aldol Condensation and Diels--Alder Reaction:\\[1mm] Tetraphenylcyclopentadienone and Its Cycloaddition with Styrene}{Organic Experiment 2}{8 October 2026}

\noindent\textbf{Abstract.} 2,3,4,5-Tetraphenylcyclopentadienone (tetracyclone, \textbf{1}) was generated in situ by the base-catalysed double aldol condensation of benzil with 1,3-diphenylacetone in triethylene glycol, and was trapped in the same vessel by a Diels--Alder reaction with styrene at 110\,\si{\celsius}. The recrystallised product was obtained as a \res{white} solid (\res{0.XXX} g, \res{XX}\%) with mp \res{XXX--XXX}\,\si{\celsius}. On the basis of IR and \textsuperscript{1}H NMR data and the X-ray crystal structure, it was identified as \res{3A / 3B / 4}.

\noindent\rule{\linewidth}{0.6pt}

% =====================================================================
\section{Introduction}

The aldol condensation and the Diels--Alder reaction are two of the most widely used carbon--carbon bond-forming reactions. Both are atom-economical,\cite{Trost1991} and in this experiment they are combined in one pot: the intermediate diene is not isolated but reacts directly with the dienophile. This saves the solvent and materials needed to isolate the intermediate, reduces the number of vessels and allows a reactive intermediate to be trapped in situ.\cite{L2}

\textbf{Aldol condensation.} In the presence of KOH, 1,3-diphenylacetone is deprotonated at a carbon that is activated by both the carbonyl group and a phenyl ring. The resulting enolate adds to one carbonyl group of benzil, and the aldol adduct dehydrates to an enone. A second, intramolecular aldol addition followed by dehydration closes the five-membered ring and gives tetracyclone \textbf{1} (Figure~\ref{fig:scheme}), a classic preparation.\cite{Johnson1943} The loss of two molecules of water drives the reaction forward. Compound \textbf{1} is deep purple because of its extended conjugation, which makes its formation easy to follow by eye and by TLC.

\begin{figure}[H]
\centering
\setchemfig{atom sep=1.35em}
\schemestart
\chemfig{Ph-[:30](=[:90]O)-[:-30](=[:-90]O)-[:30]Ph}
\+
\chemfig{Ph-[:-30]-[:30](=[:90]O)-[:-30]-[:30]Ph}
\arrow{->[\footnotesize KOH, EtOH][\footnotesize TEG, 110\,\si{\celsius}]}[,1.7]
\chemfig{*5((=O)-(-Ph)=(-Ph)-(-Ph)=(-Ph)-)}
\arrow{->[\footnotesize styrene \textbf{2}][\footnotesize 110\,\si{\celsius}]}[,1.4]
{\small \textbf{3A} (endo), \textbf{3B} (exo) or \textbf{4}}
\schemestop
\caption{One-pot sequence studied. Tetracyclone \textbf{1} (MW 384.48) reacts with styrene \textbf{2} to give the bridged ketone adducts \textbf{3A}/\textbf{3B} (1,2,3,4,5-pentaphenylbicyclo[2.2.1]hept-2-en-7-one, MW 488.63), which could also lose CO to give the cyclohexadiene \textbf{4} (1,2,3,4,5-pentaphenylcyclohexa-1,3-diene, MW 460.62). Structures as drawn in the laboratory manual.\cite{Manual2026}}
\label{fig:scheme}
\end{figure}

\textbf{Diels--Alder reaction.} The Diels--Alder reaction, first reported in 1928,\cite{Diels1928} is a concerted [4+2] cycloaddition between a conjugated diene and an alkene (the dienophile). Two new C--C $\sigma$ bonds are formed at the expense of two $\pi$ bonds, so the reaction is thermodynamically favourable. Because all bonds form in a single step without rotation, the reaction is stereospecific.\cite{L2,Clayden2012} Tetracyclone is a good diene because its diene unit is locked in the required s-\emph{cis} geometry within the ring. Styrene is a modest dienophile: it has no strong electron-withdrawing group, but its phenyl substituent activates the double bond sufficiently at 110\,\si{\celsius}.\cite{L2}

With a cyclic diene and an unsymmetrical dienophile, two diastereomeric bicyclic adducts can form. In the \emph{endo} adduct \textbf{3A}, the phenyl group from styrene points towards the longer (two-carbon) bridge; in the \emph{exo} adduct \textbf{3B}, it points away from it. Under kinetic control (lower temperature), the endo adduct usually forms faster, whereas the exo adduct is often the more stable, thermodynamic product and can dominate at higher temperature.\cite{L2} In addition, adducts of tetracyclone contain a strained bridging carbonyl group that can be lost as carbon monoxide. This happens, for example, with benzyne and with diphenylacetylene, where CO extrusion gives aromatic products.\cite{L2,Fieser1966} With styrene, loss of CO from \textbf{3A}/\textbf{3B} would give the cyclohexadiene \textbf{4}. Each of the three candidates would be formed as a racemate.

\textbf{Structure determination.} All three candidates contain three non-aromatic hydrogens (one CH and one diastereotopic \ce{CH2}), each expected as a doublet of doublets. The geminal coupling is large (11--14 Hz), while the vicinal couplings depend on the H--C--C--H dihedral angle according to the Karplus relationship.\cite{Karplus1959,Manual2026} IR spectroscopy indicates whether the bridging ketone is still present. The stereochemistry, however, is established unequivocally by the X-ray crystal structure, as in the original description of this experiment.\cite{Hoang2015}

\textbf{Solvent.} Triethylene glycol (TEG) has a very low vapour pressure and a high boiling point (287\,\si{\celsius}), tolerates sustained heating above 110\,\si{\celsius} and has low toxicity, which makes it suitable for running both steps in one heated vial.\cite{L2}

The aims of this experiment were (i) to carry out the one-pot aldol condensation and Diels--Alder reaction, monitoring each step by TLC, (ii) to purify the product by recrystallisation, and (iii) to determine which of \textbf{3A}, \textbf{3B} and \textbf{4} was formed, using the melting point, IR and \textsuperscript{1}H NMR data and the X-ray crystal structure.

% =====================================================================
\section{Experimental}

\subsection{General}
Benzil, 1,3-diphenylacetone, triethylene glycol, styrene (stabilised with 4-\emph{tert}-butylcatechol), freshly prepared 2 M ethanolic KOH, ethanol, ethyl acetate and hexane were used as supplied. TLC was performed with 10\% ethyl acetate in hexane as the eluent, and plates were visualised under UV light. The melting point was measured on a Mel-Temp apparatus. IR spectra were recorded on \res{instrument} with an ATR accessory. \textsuperscript{1}H NMR spectra were recorded at \res{400} MHz in \ce{CDCl3}; chemical shifts are reported relative to residual \ce{CHCl3} ($\delta$ 7.26).\cite{Fulmer2010}

\subsection{One-pot synthesis}
Benzil (\res{0.210} g, \res{1.00} mmol) and 1,3-diphenylacetone (\res{0.210} g, \res{1.00} mmol) were placed in a 20 mL reaction vial with a small stirrer bar, and triethylene glycol (0.50 mL) was added. The vial was heated on a heating block at 110\,\si{\celsius} until the solids had melted or dissolved to give a bright yellow solution. Ethanolic KOH (2 M, \res{5--6} drops) was added; the mixture turned \res{purple} and \res{solidified}, and heating was continued for 3--5 min. A small sample (R1) was dissolved in ethyl acetate and analysed by TLC against the two starting materials, which showed \res{complete} consumption of the starting materials.

Styrene (0.40 mL, 3.48 mmol, 3.5 equiv) was added, and the mixture was heated at 110\,\si{\celsius} for 30 min, during which its colour changed from dark purple to \res{tan}. TLC of a sample (R2) against R1 showed \res{no 1 left}.
% If tetracyclone remained: "Further styrene (1 equiv, approx. 0.11 mL) was added and heating was continued for 20 min."

\subsection{Work-up and recrystallisation}
The mixture was cooled to room temperature, ethanol (5 mL) was added and the mixture was stirred with a spatula and cooled in an ice bath. The precipitate was collected by suction filtration on a Hirsch funnel, broken up with a spatula and washed well with ethanol to remove residual styrene. The crude solid was heated to a gentle boil in a mixture of ethanol (10 mL) and ethyl acetate (10 mL), and further hot ethyl acetate (\res{X} mL) was added dropwise until the solid had dissolved. The solution was cooled to room temperature over 15--20 min and then in an ice bath. The crystals were collected by suction filtration, rinsed with a minimal volume of ethanol and dried on the filter to give the product as \res{appearance} (\res{0.XXX} g, \res{X.XXX} mmol, \res{XX}\%).
% If crystals did not form on the first day: "The solution was left to crystallise in the fume hood until the next session."

\noindent\textbf{Characterisation data.} Mp \res{XXX--XXX}\,\si{\celsius}. IR (ATR) $\tilde{\nu}_\text{max}$/\si{\per\centi\metre}: \res{values}. \textsuperscript{1}H NMR (\res{400} MHz, \ce{CDCl3}) $\delta$: \res{values}.

% =====================================================================
\section{Results}

\begin{table}[H]
\centering
\caption{Quantities of reagents (planned values from the laboratory manual\cite{Manual2026} and actual values used).}
\small
\begin{tabular}{lcccccc}
\toprule
Reagent & $M$ / g\,mol$^{-1}$ & Density & Planned & Actual & $n$ / mmol & Equiv. \\
\midrule
Benzil & 210.23 & --- & 0.21 g & \res{} & \res{1.00} & 1.00 \\
1,3-Diphenylacetone & 210.27 & --- & 0.21 g & \res{} & \res{1.00} & \res{1.00} \\
Triethylene glycol & 150.17 & 1.113 & 0.50 mL & \res{} & solvent & --- \\
KOH (2 M in EtOH) & 56.10 & --- & 5--6 drops & \res{} drops & cat. & --- \\
Styrene & 104.15 & 0.906 & 0.40 mL & \res{} & 3.48 & \res{3.5} \\
\bottomrule
\end{tabular}
\end{table}

\begin{table}[H]
\centering
\caption{Yield. Keep the row that matches the structure identified.}
\small
\begin{tabular}{lcc}
\toprule
 & If product is \textbf{3A}/\textbf{3B} (MW 488.63) & If product is \textbf{4} (MW 460.62) \\
\midrule
Theoretical yield (from 1.00 mmol benzil) & 0.488 g & 0.460 g \\
Mass of recrystallised product & \multicolumn{2}{c}{\res{0.XXX} g} \\
Isolated yield & \res{XX}\% & \res{XX}\% \\
\bottomrule
\end{tabular}
\end{table}

\begin{table}[H]
\centering
\caption{TLC data (silica, 10\% EtOAc/hexane, UV visualisation; solvent front \res{X.X} cm).}
\small
\begin{tabular}{lccl}
\toprule
Sample & Distance / cm & $R_f$ & Comment \\
\midrule
Benzil & \res{} & \res{} & \\
1,3-Diphenylacetone & \res{} & \res{} & \\
R1 (after aldol step) & \res{} & \res{} & \obs{purple spot = 1?} \\
R2 (after Diels--Alder step) & \res{} & \res{} & \obs{is the spot of 1 gone?} \\
\bottomrule
\end{tabular}
\end{table}

\begin{table}[H]
\centering
\caption{Observations.}
\small
\begin{tabularx}{\linewidth}{lX}
\toprule
Stage & Observation \\
\midrule
Solids in TEG at 110\,\si{\celsius} & \obs{melting/dissolution, bright yellow?} \\
Addition of KOH & \obs{colour change, solidification, number of drops} \\
Addition of styrene, 30 min & \obs{purple to tan? time taken} \\
Ethanol addition and cooling & \obs{precipitate} \\
Recrystallisation & \obs{volume of EtOAc added, crystal form} \\
Melting point & \obs{range; any gas evolution or colour change on heating} \\
\bottomrule
\end{tabularx}
\end{table}

\begin{table}[H]
\centering
\caption{Diagnostic spectroscopic data of the product.}
\small
\begin{tabularx}{\linewidth}{lX}
\toprule
Measurement & Result \\
\midrule
IR, 1650--1850 \si{\per\centi\metre} region & \res{C=O band at X cm$^{-1}$ / no C=O band} \\
\textsuperscript{1}H NMR, aromatic region & \res{$\delta$ range}, integral \res{X} (25H expected for all candidates) \\
\textsuperscript{1}H NMR, H$_a$ & \res{$\delta$}, dd, $J$ = \res{X}, \res{X} Hz, 1H \\
\textsuperscript{1}H NMR, H$_b$ & \res{$\delta$}, dd, $J$ = \res{X}, \res{X} Hz, 1H \\
\textsuperscript{1}H NMR, H$_c$ & \res{$\delta$}, dd, $J$ = \res{X}, \res{X} Hz, 1H \\
X-ray crystal structure (provided) & \res{3A / 3B / 4} \\
\bottomrule
\end{tabularx}
\end{table}

% =====================================================================
\section{Discussion}

\subsection{The aldol step}
The rapid change from bright yellow to deep purple on addition of KOH shows the formation of the highly conjugated tetracyclone \textbf{1}. TLC of R1 confirmed \res{R1 TLC result}. Only a catalytic amount of base is needed because hydroxide is regenerated in each dehydration step. Heating in TEG keeps the reactants molten or dissolved without solvent loss, and the product precipitates as the reaction proceeds, which helps to drive it to completion.

\subsection{The Diels--Alder step}
The fading of the purple colour to \res{tan} during heating with styrene indicates that the conjugated cyclopentadienone chromophore of \textbf{1} was consumed. In the adducts \textbf{3A}/\textbf{3B}, the former diene carbons become sp$^3$ bridgeheads, so the conjugation is broken. Styrene was used in excess (3.5 equiv) because it is only a moderately reactive dienophile and some may evaporate or polymerise at 110\,\si{\celsius}; the excess was removed by washing the crude solid with ethanol.

\subsection{Identification of the product}
\textbf{Bridging ketone or not.} Compounds \textbf{3A} and \textbf{3B} contain a carbonyl group bridging a bicyclo[2.2.1] framework. This carbonyl is strongly strained, and ring strain shifts C=O stretching bands to higher wavenumber (for comparison, cyclopentanone absorbs at a higher wavenumber than cyclohexanone).\cite{Pavia} A C=O band well above that of an unstrained ketone is therefore expected for \textbf{3A}/\textbf{3B}, whereas \textbf{4} has no carbonyl group at all. The IR spectrum showed \res{C=O at X cm$^{-1}$ / none}, which \res{supports 3A/3B / supports 4}.

\textbf{Non-aromatic protons.} All three candidates have 25 aromatic and 3 non-aromatic protons, so the integral ratio (\res{25.0 : 3.0}) confirms that a 1:1 adduct was formed but cannot distinguish between them. The three non-aromatic signals appeared as doublets of doublets. The largest coupling, \res{X} Hz, lies in the 11--14 Hz range expected for the geminal pair of \ce{CH2} protons.\cite{Manual2026} The vicinal couplings between the benzylic CH and the two \ce{CH2} protons follow the Karplus relationship:\cite{Karplus1959} couplings are largest for dihedral angles near 0\si{\degree} or 180\si{\degree} and smallest near 90\si{\degree}. In the rigid bicyclic skeleton of \textbf{3A}/\textbf{3B}, the hydrogens are held in eclipsed positions, so the benzylic H is \emph{cis} to one \ce{CH2} proton (dihedral angle close to 0\si{\degree}, large $J$) and \emph{trans} to the other (larger dihedral angle, smaller $J$). In the flexible cyclohexadiene ring of \textbf{4}, the dihedral angles are instead averaged over ring conformations. \resp{Interpret your two vicinal J values here (which CH$_2$ proton is cis and which is trans to the benzylic H).}
% Note: both 3A and 3B have one cis and one trans CH2 proton relative to the benzylic H, so J values alone may not decide endo vs exo; this is why the X-ray structure is needed.

\textbf{X-ray structure.} The crystal structure provided showed \res{3A / 3B / 4}, in \res{agreement} with the spectroscopic interpretation. This illustrates that IR and NMR spectroscopy can narrow down the possibilities, but an X-ray structure is needed to establish the relative stereochemistry unequivocally.\cite{Hoang2015}
% If 3A (endo): consistent with kinetic control; the endo transition state benefits from secondary orbital interactions.
% If 3B (exo): at 110 C the reaction may be under (partial) thermodynamic control, or steric repulsion with the four phenyl groups of 1 may disfavour the endo approach.
% If 4: the temperature was high enough for cheletropic loss of CO from the initial adduct.

\textbf{Melting point.} The product melted at \res{XXX--XXX}\,\si{\celsius}. The laboratory manual notes that the product may release styrene when heated above 170\,\si{\celsius},\cite{Manual2026} that is, it may undergo a retro-Diels--Alder reaction. \resp{Comment on any gas evolution or colour change observed near the melting point.}

\subsection{Yield and efficiency}
The isolated yield was \res{XX}\%. \resp{Discuss losses: material left on the vial and Hirsch funnel, solubility in the ethanol washes and in the recrystallisation mother liquor, incomplete reaction (TLC of R2).}

The sequence is highly atom-economical. The aldol step loses only two molecules of water, giving an atom economy of $384.48/(210.23+210.27) = 91.4\%$ for \textbf{1}; the Diels--Alder step itself incorporates all atoms of both partners. The overall atom economy from benzil, 1,3-diphenylacetone and styrene is 93.1\% for \textbf{3A}/\textbf{3B}, or 87.8\% for \textbf{4} if CO is lost.\cite{Trost1991} Running both steps in one vessel also avoided isolating, drying and redissolving tetracyclone.

% =====================================================================
\section{Conclusion}
Tetracyclone was prepared by a double aldol condensation and trapped in situ with styrene in a one-pot sequence in triethylene glycol. The product was obtained in \res{XX}\% yield after recrystallisation and identified as \res{3A / 3B / 4} on the basis of its IR and \textsuperscript{1}H NMR spectra and the X-ray crystal structure. The experiment shows how a one-pot sequence of two atom-economical reactions can build a complex, polyaryl bicyclic framework, and why X-ray crystallography is needed to confirm stereochemistry that spectroscopy alone can only suggest.

% =====================================================================
\begin{thebibliography}{99}
\small
\setlength{\itemsep}{1pt}
\bibitem{L2} Hoang, T.\,G. \emph{CM3191 Org L2: One-Pot Aldol Condensation and Diels--Alder Reaction} (lecture slides); National University of Singapore, 2026.
\bibitem{Johnson1943} Johnson, J.\,R.; Grummitt, O. Tetraphenylcyclopentadienone. \emph{Org. Synth.} \textbf{1943}, \emph{23}, 92.
\bibitem{Manual2026} Hoang, T.\,G. \emph{CM3191 Chemical Synthesis Experiments: Organic Lab Manual}; National University of Singapore, AY2026/27; Experiment 2, pp 8--13.
\bibitem{Diels1928} Diels, O.; Alder, K. Synthesen in der hydroaromatischen Reihe. \emph{Justus Liebigs Ann. Chem.} \textbf{1928}, \emph{460}, 98--122.
\bibitem{Clayden2012} Clayden, J.; Greeves, N.; Warren, S. \emph{Organic Chemistry}, 2nd ed.; Oxford University Press: Oxford, 2012.
\bibitem{Fieser1966} Fieser, L.\,F. Hexaphenylbenzene. \emph{Org. Synth.} \textbf{1966}, \emph{46}, 44.
\bibitem{Karplus1959} Karplus, M. Contact Electron-Spin Coupling of Nuclear Magnetic Moments. \emph{J. Chem. Phys.} \textbf{1959}, \emph{30}, 11--15.
\bibitem{Hoang2015} Hoang, G.\,T.; Kubo, T.; Young, V.\,G.; Kautzky, J.\,A.; Wissinger, J.\,E. Illustrating the Utility of X-ray Crystallography for Structure Elucidation through a Tandem Aldol Condensation/Diels--Alder Reaction Sequence. \emph{J. Chem. Educ.} \textbf{2015}, \emph{92}, 1381--1384.
\bibitem{Fulmer2010} Fulmer, G.\,R.; Miller, A.\,J.\,M.; Sherden, N.\,H.; Gottlieb, H.\,E.; Nudelman, A.; Stoltz, B.\,M.; Bercaw, J.\,E.; Goldberg, K.\,I. NMR Chemical Shifts of Trace Impurities: Common Laboratory Solvents, Organics, and Gases in Deuterated Solvents Relevant to the Organometallic Chemist. \emph{Organometallics} \textbf{2010}, \emph{29}, 2176--2179.
\bibitem{Pavia} Pavia, D.\,L.; Lampman, G.\,M.; Kriz, G.\,S.; Vyvyan, J.\,R. \emph{Introduction to Spectroscopy}, 5th ed.; Cengage Learning: Stamford, CT, 2015.
\bibitem{Trost1991} Trost, B.\,M. The Atom Economy---A Search for Synthetic Efficiency. \emph{Science} \textbf{1991}, \emph{254}, 1471--1477.
\end{thebibliography}

\end{document}
